\section{Introduction}
\label{sec:introduction}

A routine task in molecular biology is to characterise a nucleotide or protein
sequence and then look at the three-dimensional structure of the protein it
encodes. In practice this is a chain of separate tools: a sequence-analysis
package for open reading frames and composition, a primer-design page, a protein
property calculator, and then one or more structure resources, namely the
experimental archive RCSB PDB \citep{berman2000pdb}, the predicted-model
database AlphaFold DB \citep{varadi2022alphafolddb,jumper2021alphafold}, and the
language-model predictor ESM Atlas \citep{lin2023esmfold}. Each hand-off is a
manual copy-and-paste, and each is a place where an identifier can be
mistyped or a result misread.

The step that is most easily skipped is the one that matters most for the
conclusion drawn from a structure: confirming that the structure that came back
is the protein that was asked about. A sequence search is designed to return
\emph{similar} sequences, and the structure archives are keyed in different
ways: AlphaFold DB by UniProt accession, RCSB PDB by entry and by sequence
similarity, ESM Atlas by predicting from the sequence alone. A hit therefore
tells the user that something was found, not that it is the queried molecule: it
may be a paralogue, a fragment, a different isoform, or the same protein
catalogued under a different accession. Scripts that fetch ``the first
structure'' and report success make this distinction invisible, and the question
of how often it matters has, to our knowledge, rarely been measured.

This paper has two parts. The first is a software contribution. \bsi\ 2.0
integrates sequence analysis with structure retrieval and, for every retrieved
structure, performs the check that is usually skipped: the structure's chain is
aligned to the query, the structure's own database cross-references are compared
with the queried accession, and the outcome is reported on an explicit five-level
scale (M0--M4) in which only a full-length structure whose accession is confirmed
is called an exact match (Section~\ref{sec:m0m4}). The second is a measurement.
Using \N{ds-n} randomly sampled reviewed UniProt proteins from \N{ds-orgs}
organisms, we ask how often accepting the top sequence-search hit returns a
structure that is not catalogued under the queried accession, how often the
validation layer prevents that, what it costs in correct structures rejected,
and where its thresholds are and are not meaningful
(Section~\ref{sec:validation-study}). Correctness is judged against curated
cross-references from PDBe SIFTS and RCSB rather than against the software's own
alignment, so the evaluation does not simply restate the classifier's output.

A note on how this study came to be framed, because it changed the paper. An
earlier prototype of this software (v1.0) reported a structural hit for every
query in its small benchmark, and its own report flagged that for 7 of 25
queries the UniProt accession attached to the hit differed from the query.
v1.0 had no step that compared the retrieved structure with the query sequence,
so it could not say whether those hits were wrong. We initially read the flags
as failures. Re-adjudicating all 25 queries against current data
(Section~\ref{sec:case-study}) shows that for \N{cs-disc-contains} of the
\N{cs-disc-n} flagged queries the top RCSB sequence-search hit is an entry that
does contain the queried accession (\N{cs-disc-M4} validate as M4), so the
discordance reflected how accessions were attributed to hits, not wrong structures. The lesson we take is methodological:
an accession mismatch in a reporting layer is not the same thing as a wrong
structure, and the two should be separated by an explicit sequence and
cross-reference check on a large sample, which is what the present study does.

\paragraph{Contributions.}
(i) An integrated, dependency-light platform with a command-line and a graphical
interface (Sections~\ref{sec:architecture}--\ref{sec:export}).
(ii) A five-level sequence-to-structure classification with a partial-structure
qualifier, applied uniformly to three heterogeneous resources
(Section~\ref{sec:m0m4}).
(iii) A measurement on \N{ds-n} random proteins of how often naive retrieval
returns a structure not catalogued under the queried accession, how often
validation prevents it, and what that costs
(Section~\ref{sec:validation-study}), with an adversarial benchmark of
constructed negatives (Section~\ref{sec:adversarial}).
(iv) Engineering evidence: an offline test suite, deterministic fault
injection, cross-validation against an established implementation, performance
measurements and an interaction-cost comparison
(Sections~\ref{sec:testing}--\ref{sec:workflow-comparison}).
